% Generated draft: requires mathematical and source review.
\documentclass[11pt]{article}
\usepackage[margin=25mm]{geometry}
\usepackage{fontspec}
\setmainfont{Latin Modern Roman}
\usepackage{amsmath,amssymb,mathtools,unicode-math}
\setmathfont{Latin Modern Math}
\usepackage{enumitem,needspace,xurl,hyperref}
\setlength{\parindent}{0pt}
\setlength{\parskip}{5pt}
\setlength{\emergencystretch}{3em}
\newcommand{\N}{\mathbb{N}}
\newcommand{\Z}{\mathbb{Z}}
\newcommand{\Q}{\mathbb{Q}}
\newcommand{\R}{\mathbb{R}}
\newcommand{\C}{\mathbb{C}}
\newcommand{\F}{\mathbb{F}}
\newcommand{\A}{\mathbb{A}}
\newcommand{\PP}{\mathbb{P}}
\newcommand{\E}{\mathbb{E}}
\newcommand{\eps}{\varepsilon}
\newcommand{\dd}{\,\mathrm{d}}
\begin{document}
\section*{UnsolvedMath ID: 701197}
\textbf{Source problem code:} catalogue entry
\subsection{Definitions and mathematical statement}
Let \(p\) be a prime. A group \(G\) is a \(p\)-group if every element has finite order equal to a power of \(p\), and is locally finite if every finitely generated subgroup is finite. For a group \(H\), write \(H'=[H,H]\) for its derived subgroup. The group \(G\) is perfect if \(G=G'\).

Define the upper central series of \(H\) by \(Z_0(H)=1\), by letting \(Z_{\alpha+1}(H)\) be the inverse image of \(Z(H/Z_\alpha(H))\) under the quotient map, and by setting \(Z_\lambda(H)=\bigcup_{\alpha<\lambda}Z_\alpha(H)\) for every limit ordinal \(\lambda\). The group \(H\) is hypercentral if \(Z_\alpha(H)=H\) for some ordinal \(\alpha\). Define \(H^{(0)}=H\) and \(H^{(n+1)}=[H^{(n)},H^{(n)}]\) for integers \(n\geq 0\). The group \(H\) is solvable if \(H^{(n)}=1\) for some finite integer \(n\geq 0\).

The two existence questions proposed by A. O. Asar are:
\begin{enumerate}
\item[(a)] Does there exist a perfect locally finite \(p\)-group \(G\) such that every proper subgroup \(H<G\) is hypercentral?
\item[(b)] Does there exist a perfect locally finite \(p\)-group \(G\) such that every proper subgroup \(H<G\) is solvable?
\end{enumerate}
These are separate existence questions; they do not require the same group to witness both properties. The ordinal in (a) and the finite derived length in (b) may depend on \(H\); no uniform bound is stated.

\emph{Unresolved qualifications.} The supplied statement does not explicitly require \(G\) to be nontrivial, although the definition of perfectness permits the trivial group. This possible missing restriction requires review and has not been added. The source also does not specify whether existence is requested for some prime \(p\), for every prime \(p\), or for a particular fixed prime. The questions above retain \(p\) as a prime parameter without resolving that quantifier.
\subsection{Short English statement}
A. O. Asar asks separately whether a perfect locally finite \(p\)-group can have every proper subgroup hypercentral, and whether such a group can have every proper subgroup solvable. The intended restriction on trivial groups and the quantification over primes require clarification. The supplied registry labels the problem open; this label is not independently verified.
\subsection{Sources}
\begin{itemize}
\item Local UnsolvedMath registry, permanent ID \(701197\), exact source problem code \texttt{catalogue entry}, title \emph{Kourovka Notebook Problem 16.5}. Supplied upstream revision: \texttt{d371ee917ac9647fbb06cd6fc529deec6f3dda82}.
\item Supplied attribution: A. O. Asar, proposed in 2006; original problem 16.5 in \emph{The Kourovka Notebook (Unsolved Problems in Group Theory)}, 16th issue (2006).
\item The supplied background also identifies the 21st issue edition (2026), \texttt{arXiv:1401.0300v48}. The supplied collection description names E. I. Khukhro and V. D. Mazurov as editors. These bibliographic details are supplied source information and have not been independently checked.
\item The registry records status \texttt{open} and an update timestamp of \texttt{2026-10-16T14:00:00Z}. This timestamp is registry metadata, not a supplied literature-verification date. No separate literature assessment or external reference links are supplied.
\end{itemize}
\end{document}
